\section{Exact support for quadratic blocks}\label{sec:quadratic}

If every feature $f_j$ is a quadratic polynomial, then every support
problem~\eqref{eq:support} asks for the minimum of a single quadratic
function over the block domain. This section gives two exact rational
algorithms for that problem. The first works for any bounded polytope of
small dimension; the second works for quadratic stars of any size whose
constraints couple the center with one leaf at a time. Both return an
attained minimizer and a rational minimum, so they serve as exact
support oracles and give support inequalities that are tight.

\subsection{Bounded polytopes in fixed dimension}\label{sec:polytope}

Let $P=\{x\in\R^d:Ax\le b\}$ be a bounded polytope given by $m$ rational
rows, which include finite variable bounds, and let
$q(x)=\tfrac12x\T Hx+c\T x+c_0$ with rational data and $H=H\T$. The
matrix $H$ may be indefinite or singular, $P$ may be lower-dimensional,
and rows may be redundant or implied equalities; an equality is given by
two opposite rows. For a set $S$ of rows, $A_S$ and $b_S$ denote the
corresponding submatrix and subvector, and
\begin{equation}\label{eq:kkt-system}
K_S=\begin{pmatrix}H&A_S\T\\ A_S&0\end{pmatrix},\qquad
K_S\begin{pmatrix}x\\ \lambda\end{pmatrix}=\begin{pmatrix}-c\\ b_S\end{pmatrix}.
\end{equation}
The candidates are the points $x$ obtained from all sets $S$ with
$\abs S\le d$ for which $K_S$ is nonsingular and $Ax\le b$ holds. The
multipliers $\lambda$ are not sign-restricted: the system expresses
stationarity on the affine hull of a face, not a KKT condition.

\begin{theorem}[Exact quadratic support on a polytope]\label{thm:polytope}
If $P\ne\emptyset$, some candidate minimizes $q$ over $P$; hence the
least candidate value is $\min_Pq$, and it is attained at a rational
point. If there is no candidate, then $P=\emptyset$. Both statements
remain true if only sets $S$ of linearly independent rows are enumerated.
\end{theorem}

\begin{proof}
Every candidate lies in $P$, so an empty $P$ has none. Assume
$P\ne\emptyset$. Among the global minimizers, which exist by
compactness, choose $x^\ast$ such that the face $F$ of $P$ with
$x^\ast\in\relint F$ has the smallest dimension, and let $T$ be the
linear space parallel to $\aff F$. Because $x^\ast$ minimizes $q$ over
$F$ and lies in its relative interior,
$v\T(Hx^\ast+c)=0$ and $v\T Hv\ge0$ for all $v\in T$.

Suppose $v\in T\setminus\{0\}$ satisfies $v\T Hv=0$. Then
$q(x^\ast+tv)=q(x^\ast)$ for all $t$. The line $x^\ast+tv$ meets the
bounded face $F$ in a segment that contains $x^\ast$ in its relative
interior, and an endpoint of the segment is a global minimizer on a
proper face of $F$, which contradicts the choice of $x^\ast$. Hence $H$
is positive definite on $T$, or $T=\{0\}$.

Let $I$ be the set of rows active at $x^\ast$, so that
$\aff F=\{x:A_Ix=b_I\}$ and $T=\ker A_I$, and let $S\subseteq I$ index
a basis of the row space of $A_I$. Then $\abs S\le d$ and
$\ker A_S=T$. Since $Hx^\ast+c$ is orthogonal to $T$, it lies in the row
space of $A_S$, so $(x^\ast,\lambda)$ solves~\eqref{eq:kkt-system} for
some $\lambda$. If $K_S(v,\mu)=0$, then $A_Sv=0$, so $v\in T$, and
$v\T Hv=-v\T A_S\T\mu=0$; hence $v=0$, and $A_S\T\mu=0$ gives $\mu=0$
by independence. So $K_S$ is nonsingular, its unique solution contains
$x^\ast$, and $x^\ast$ is a candidate. All data are rational, so the
candidates and their values are rational.
\end{proof}

The sign of the multipliers and the inertia of $H$ on $T$ need not be
tested. Extra candidates, such as saddle points of faces, are feasible
and therefore cannot produce a value below the minimum. A set $S$ with
dependent rows never gives a nonsingular $K_S$, which is why independent
sets suffice; duplicate, redundant and opposite rows are therefore
harmless. Two parts of the argument cannot be dropped. Boundedness is
needed: for $P=\{x\in\R:x\ge0\}$ and $q=-x$ the enumeration returns
$0$, and for $P=\R$ and $q=x$ it returns no candidate. The choice of a
minimizer on a face of minimal dimension is also needed: for
$q=(x-y)^2$ on $[0,1]^2$, the minimizers $(t,t)$ with $0<t<1$ give only
singular systems, and the enumeration finds the minimizers $(0,0)$ and
$(1,1)$ instead.

\begin{corollary}[Support of a quadratic feature graph]\label{cor:quad-features}
If $f_1,\dots,f_k$ are quadratic polynomials with rational coefficients
and $(a,b)$ is rational, the support value $h_P(a,b)$ of
$\conv\{(x,f(x)):x\in P\}$ is the rational number obtained by applying
Theorem~\ref{thm:polytope} to $q=a\T x+b\T f$, and the support
inequality is tight at the returned minimizer.
\end{corollary}

There are at most $N(m,d)=\sum_{k=0}^d\binom mk$ sets $S$, where $m$
counts the box rows as well. For each, Gaussian elimination on a system
of order at most $2d$ and the feasibility test take $O(d^3+md)$
arithmetic operations, so the enumeration uses $O(N(m,d)(d^3+md))$
operations. With Bareiss fraction-free elimination, every intermediate
number is a minor of an integer matrix of order at most $2d$, and
Hadamard's inequality bounds the bit length of every candidate and of
the minimum by $O(d(\lambda_q+\lambda_A+\log d))$, where $\lambda_q$ is
the encoding length of $q$ and $\lambda_A$ the largest encoding length of
a row (Appendix~\ref{app:polytope}). These sizes are polynomial for every
$d$; only the number of systems grows exponentially with $d$. For fixed
$d$ the algorithm is therefore strongly polynomial, and the minimum has
polynomial encoding length, in line with the result that quadratic
programming is in NP \citep{Vavasis1990}. The method itself is total
enumeration of faces \citep[Section~2.9]{Murty1997}; we give the proof
because the certificate and its replay depend on the exact treatment of
the degenerate cases above.

\begin{proposition}[Hardness in unrestricted dimension]\label{prop:maxcut}
Computing $\min_{x\in[0,1]^n}q(x)$ for rational quadratics $q$ is strongly
NP-hard, even when $q$ is multilinear with integer coefficients of
absolute value at most $n$. Deciding whether a given inequality
$q(x)\ge t$ holds on $[0,1]^n$ is coNP-complete.
\end{proposition}

\begin{proof}
For a graph $G=(V,E)$ let $q(x)=-\sum_{\{i,j\}\in E}(x_i+x_j-2x_ix_j)$.
Since $q$ is affine in each coordinate, moving coordinates one at a time
to a minimizing endpoint does not increase $q$, so some minimizer is
binary. At a binary point the summand of an edge equals one exactly
when the edge crosses the cut defined by $x$, so $\min q$ is minus the
maximum cut size, and unweighted MAX CUT is NP-hard
\citep{Karp1972,GareyJohnsonStockmeyer1976}. For the second statement,
membership in coNP follows because a violating point can be taken to be
a candidate of Theorem~\ref{thm:polytope}, which has polynomial size; with
integer data, $\min q\le-W$ holds if and only if $q\ge-W+\tfrac12$ fails.
\end{proof}

The proposition applies to the box itself, with no further rows. Joint
support for blocks of unrestricted size is therefore as hard as
nonconvex quadratic programming, and, unless NP${}={}$coNP, the validity
of an exact support cut in unrestricted dimension has no short certificate
that can be checked in polynomial time. A certificate that is replayed by
repeating the enumeration, as ours is, is thus natural for this class.
Useful exact oracles must restrict the dimension, as in
Theorem~\ref{thm:polytope}, or the structure, as in
Section~\ref{sec:star}. Restricting the interaction graph alone is not
enough in general: box-constrained quartic minimization is strongly
NP-hard already on paths \citep{DelPiaKhajavirad2026}.

\subsection{Why the domain of the block matters}\label{sec:domain-matters}

Affine rows that involve only block variables belong in the support
domain $D$, not only in the LP. On the triangle
$T=\{(x,y):x,y\ge0,\ x+y\le1\}$ every point satisfies
\[
x^2+2xy+y^2\le x+y\qquad\text{and}\qquad xy\le\tfrac14,
\]
because $s=x+y\in[0,1]$ gives $s^2\le s$ and $4xy\le s^2$. In the lifted
coordinates $(x,y,X,Y,Z)$ for $(x,y,x^2,y^2,xy)$ these are linear
inequalities valid for $\conv\{(x,y,x^2,y^2,xy):(x,y)\in T\}$. The point
obtained by averaging $(0,0)$ and $(1,1)$ in the lifted space of the box
$[0,1]^2$ is $(\tfrac12,\tfrac12,\tfrac12,\tfrac12,\tfrac12)$. It lies
in the convex hull of the box graph, and its first-order coordinates
satisfy $x+y\le1$, but $X+2Z+Y=2>1=x+y$ and $Z=\tfrac12>\tfrac14$. Thus
intersecting the hull over the box with the row is strictly weaker than
convexifying over the row-constrained domain. The principle is known:
factorable relaxations ignore linking constraints, and convexifying over
the constrained domain repairs this \citep{Rikun1997,ZhuHeTawarmalani2026,WuMutsNowakHendrix2025}.
Here the first inequality is the sum of the two
reformulation--linearization products $(1-x-y)x\ge0$ and $(1-x-y)y\ge0$,
and on a triangle the semidefinite and RLT constraints together describe
the hull exactly \citep{AnstreicherBurer2010}. The example therefore
shows the weakness of the box hull, not an advantage over RLT or
semidefinite relaxations. Theorem~\ref{thm:polytope} works on the
row-constrained domain directly and applies to any polytope of small
dimension. For a single bilinear term, SCIP's projection-based envelopes
\citep{MullerSerranoGleixner2020} exploit linear rows in the same
spirit.

\subsection{Quadratic stars with center--leaf constraints}\label{sec:star}

Theorem~\ref{thm:polytope} is limited to small dimension, and
Section~\ref{sec:composition} shows that the blocks worth merging are
often stars: several pairs that share one variable. For stars there is an
exact algorithm whose cost is nearly linear in the size of the block.

Let $y$ be the center and $x_1,\dots,x_k$ the leaves, and consider
\begin{equation}\label{eq:star}
q(y,x)=a_0+a_1y+a_2y^2+\sum_{i=1}^k\bigl(d_ix_i^2+(e_iy+f_i)x_i\bigr)
\end{equation}
with rational coefficients of arbitrary signs. The domain $S$ is given by
finite rational bounds on all variables and by rational affine rows, each
of which involves $y$ and \emph{at most one} leaf; an equality is given as
two rows. Every scalarization $a\T(y,x)+b\T F(y,x)$ of a vector $F$ of
quadratics with this interaction pattern has the form~\eqref{eq:star}.
Rows or products that couple two leaves are excluded, and for good reason:
minimizing a concave separable quadratic over a box and a single row that
couples all variables is already NP-hard \citep{MoreVavasis1990}.

For fixed $y$, the rows of leaf $i$ restrict $x_i$ to an interval
$[L_i(y),U_i(y)]$, where $L_i$ is the maximum of the affine lower bounds and
$U_i$ the minimum of the affine upper bounds that these rows and the leaf
bounds define. Thus $L_i$ is convex and $U_i$ concave, both piecewise
affine. The set $J_i=\{y:L_i(y)\le U_i(y)\}$ is an interval because $L_i-U_i$
is convex. Given $y$, the leaves are constrained independently, so the
projection of $S$ onto $y$ is the interval $I$ obtained by intersecting
the bounds of $y$, the rows in $y$ alone, and all $J_i$. If $I$ is empty,
so is $S$.

\begin{theorem}[Exact support on constrained stars]\label{thm:star}
Let $m$ be the number of rows that involve a leaf, and suppose $S$ is
nonempty. Then $\min_Sq$ is attained at a rational point, and it can be
computed, together with a minimizer and a partition of $I$ into at most
$2m+4k+1$ intervals on each of which every leaf has an affine minimizer,
with $O((m+k)\log(m+k))$ arithmetic operations and comparisons. All
numbers involved have bit length polynomial in the input length; the
breakpoints have bit length bounded independently of $k$.
\end{theorem}

\begin{proof}
Define $\varphi_i(y)=\min\{d_ix_i^2+(e_iy+f_i)x_i:L_i(y)\le x_i\le U_i(y)\}$
for $y\in I$, so that $\min_Sq=\min_{y\in I}V(y)$ with
$V(y)=a_0+a_1y+a_2y^2+\sum_i\varphi_i(y)$. Let $s_i=m_i+2$ be the number of
affine functions that define $L_i$ and $U_i$, where $m_i$ is the number of
rows of leaf $i$; together $L_i$ and $U_i$ have at most $s_i-2$ breakpoints.

If $d_i>0$, the minimizer is $\xi_i=\max\{L_i,\min\{v_i,U_i\}\}$ with the
affine function $v_i(y)=-(e_iy+f_i)/(2d_i)$. The function $v_i-L_i$ is
concave and $v_i-U_i$ convex, so each has at most two isolated zeros, and
$\xi_i$ is affine between consecutive points of a set of at most $s_i+2$
breakpoints. If $d_i\le0$, a minimizer is an endpoint of the interval, and
\begin{equation}\label{eq:endpoint-factor}
\psi_i(U_i)-\psi_i(L_i)=(U_i-L_i)\bigl(d_i(U_i+L_i)+e_iy+f_i\bigr),
\qquad\psi_i(t)=d_it^2+(e_iy+f_i)t .
\end{equation}
The first factor is nonnegative on $I$. On each of the at most $s_i-1$
intervals where $L_i$ and $U_i$ are affine, the second factor is affine and
has at most one zero, which decides the preferred endpoint; this gives at
most $2s_i-3$ breakpoints. Taking all breakpoints together splits $I$ into
at most $\sum_i2s_i+1=2m+4k+1$ intervals on which every leaf rule is affine
and $V(y)=q(y,\xi(y))$ is a quadratic polynomial in $y$.

Each $\varphi_i$ is continuous on $I$, since it is the minimum of a
continuous function over an interval whose endpoints depend continuously
on $y$. Hence $V$ agrees with the quadratic of each interval on its
closure, and the minimum of $V$ is the least value of these quadratics at
the interval endpoints and, where the leading coefficient is positive, at
interior stationary points. All breakpoints are zeros of affine functions
with rational coefficients, so every candidate is rational, and the
minimizing $y$ together with $x_i=\xi_i(y)$ is a rational minimizer.

For the operation count, the upper envelope of $s_i$ lines is computed by
sorting slopes and one stack-based scan in $O(s_i\log s_i)$ operations, and
the breakpoints and rules of leaf $i$ in $O(s_i)$ more. Sorting all
breakpoints costs $O((m+k)\log(m+k))$. A sweep over the sorted breakpoints
maintains the three coefficients of $V$; at each breakpoint only the
leaves whose rule changes are updated, so the sweep takes $O(m+k)$
operations. Bit lengths are bounded in Appendix~\ref{app:star}.
\end{proof}

The counts in the proof are attained. A single concave leaf whose bounds
zigzag forces $2s_i-3$ rule changes, and $k$ hinge leaves force $k$
more, so the partition can have $\Theta(m+k)$ intervals. On instances
with $m=\Theta(k)$, the operation bound is optimal up to a constant factor
for algebraic computation trees: deciding whether the minimum is at least
a given value requires $\Omega((m+k)\log(m+k))$ operations, by Ben-Or's
theorem applied to a family whose decision set has $r!$ components
(Appendix~\ref{app:star}).

\paragraph{Relation to forest algorithms.}
\citet{DelPiaKhajavirad2026} solve box-constrained quadratic programs on
forests with $O(n^2)$ arithmetic operations. After an affine scaling of
each variable to $[0,1]$, a box-constrained star is such a forest, and
their dynamic program rooted at the center specializes to the
box-constrained case of Theorem~\ref{thm:star}; in that case our
contribution is only the $O(k\log k)$ count, which replaces their
$O(k^2)$ root merge. Rows that couple the center with a leaf are outside
their model. A shear $x_i\mapsto x_i-\delta y$ turns a leaf domain that is a
parallelogram with two vertical sides back into a box and keeps the star
structure, but triangles, trapezoids with nonparallel sides and polygons
with more sides cannot be reduced in this way. Such rows also break the
structural property on which their algorithm rests: the conditional value
of a leaf is then no longer a minimum of affine functions of $y$ over a
fixed set, and it can have convex kinks. For example, for $x\in[0,1]$
with rows $x\ge y-\tfrac12$ and $x\ge\tfrac12-y$, the minimum of $x$ is
$\abs{y-\tfrac12}$. The SOC descriptions of star hulls by
\citet{DeyKhajavirad2025} assume sign patterns of the square terms and
box domains, and the LP approximations of \citet{BienstockMunoz2018}
allow constraints but are approximate. We know of no earlier exact
algorithm for nonconvex stars with center--leaf rows; the construction
itself, conditioning on the center and solving parametric scalar
problems, is elementary.

\begin{remark}[Deeper trees]\label{rem:trees}
The argument uses that every leaf is a leaf. On the path
$x_1$--$x_2$--$x_3$ over $[0,1]^3$ with
$f=(x_2-\tfrac18)x_1+x_2^2-x_2+x_2x_3+x_3^2-x_3$, the conditional value of the
subtree $\{x_1,x_2\}$ given $x_3=t$ is
$\min\{-\tfrac18,\,-\tfrac14(1-t)^2\}$, which switches at the irrational
point $t=1-\sqrt2/2$, although $\min f=-\tfrac38$ is rational
(Appendix~\ref{app:star}). The rational partition above therefore does
not extend by nesting conditional values. \citeauthor{DelPiaKhajavirad2026}
handle such breakpoints for boxes with an explicit algebraic
representation; we do not know whether edge-local affine rows can be
added to deeper trees with a polynomial algorithm.
\end{remark}

The certificate of a star support value records the partition of $I$,
the affine rule of every leaf on every piece, the quadratic of $V$ on
every piece and the candidates; a checker rebuilds the partition from the
rows and coefficients and compares. Our implementation is simpler than
the algorithm of the proof: it intersects all pairs of envelope lines and
recomputes every leaf rule on every piece, which costs $O((m+k)^3)$
operations and stores $\Theta(k)$ rules per piece. This is adequate for the
blocks of Section~\ref{sec:computations}, where $k\le4$.

